\section*{Chapitre \ref{ch:nw:equities-and-options-connectivity} --- Connectivité des actions et des options}

\begin{solution}{exo:nw:equities-and-options-connectivity:1}
$4.403 \times 10 = \qty{44.0}{\giga\bit\per\second}$ sur 100 millisecondes et $0.501 \times 100 = \qty{50.1}{\giga\bit\per\second}$ sur
10 millisecondes.
\end{solution}

\begin{solution}{exo:nw:equities-and-options-connectivity:2}
$2\,400\,000/50 = 48\,000$ messages par seconde, si chaque message est plein.
\end{solution}

\begin{solution}{exo:nw:equities-and-options-connectivity:3}
Chaque nuage négocie toutes les options sur son propre ensemble de sous-jacents, et une cotation de masse n'atteint que le moteur situé derrière le port par lequel elle est envoyée
sur ce port~: une firme qui cote un sous-jacent doit disposer d'un port Full Service sur le nuage de ce sous-jacent, quelle que soit la charge.
\end{solution}

\begin{solution}{exo:nw:equities-and-options-connectivity:4}
$600\,000 \times 20.4/(12 \times 2 \times 50) = 10\,200$~µs, \qty{10.2}{\milli\second}~; avec trois ports par moteur
\qty{6.8}{\milli\second}. Sur MIAX, le troisième port de cotation de masse n'est pas disponible~: les firmes obtiennent deux ports Full Service par nuage.
\end{solution}

\begin{solution}{exo:nw:equities-and-options-connectivity:5}
Sur une minute chargée de 60\,000 millisecondes, elle ne laisse que six dépassements de budget, et les millisecondes les plus chargées dépassent plusieurs fois la
moyenne sur 10 millisecondes. Au-delà, le trafic s'accumule dans les tampons du commutateur, ajoutant du délai, ou les déborde, et des paquets sont perdus, laissant des trous
que les handlers doivent combler à partir de l'autre flux ou par retransmission.
\end{solution}

\begin{solution}{exo:nw:equities-and-options-connectivity:6}
Phlx 28\,440~\$ (avant remises sur volume), EDGX 18\,000~\$, Arca et American 12\,240~\$ chacun, BOX 1\,000~\$, MIAX 20\,500~\$ pour plus de 100
classes, quel que soit le nombre de ports.
\end{solution}

\begin{solution}{exo:nw:equities-and-options-connectivity:7}
41 cœurs~: la file culmine à 199\,500 messages et se vide en \qty{0.49}{\milli\second}~; avec 40, il faut \qty{0.52}{\milli\second}.
\end{solution}

\begin{solution}{exo:nw:equities-and-options-connectivity:8}
Quatre erreurs~: les bits de l'avis sont comptés avant l'encapsulation Ethernet (\qty{59.4}{\giga\bit\per\second} sur le support)~; la marge ajoute 10~\%~;
le pic sur 10 millisecondes est une moyenne, et la milliseconde la plus chargée atteint \qty{234}{\giga\bit\per\second} dans la simulation du chapitre~;
et les lignes sont regroupées en blocs entiers, si bien que les liens ne peuvent pas être remplis uniformément. Six liens portent la moyenne sur 10 millisecondes, et non les millisecondes qu'elle contient.
\end{solution}

\begin{solution}{pb:nw:equities-and-options-connectivity:1}
\textbf{Partie I.}
\begin{enumerate}
\item Selon le décompte de l'avis, \qty{44.0}{\giga\bit\per\second} sur 100 millisecondes et \qty{50.1}{\giga\bit\per\second} sur 10~; sur le
  support 52.2 et \qty{59.4}{\giga\bit\per\second}~; les deux flux avec la marge 114.8 et \qty{130.6}{\giga\bit\per\second}.
\item \qty{48.0}{bytes}~: les 352 octets par paquet de l'avis, pris comme charge utile UDP, encapsulés à 418 octets sur le support, rapportés à 8.68 messages par paquet.
\item 4, 12 et 44.
\item Les lignes sont regroupées en blocs entiers, et chaque ligne a ses propres rafales, jusqu'à \qty{6.8}{\giga\bit\per\second} sur une milliseconde, à des
  instants différents~; si bien que presque chaque ligne demande son propre lien, contre une borne de 26 par flux pour la seule milliseconde la plus chargée.
\end{enumerate}
\textbf{Partie II.}
\begin{enumerate}[resume]
\item 16 par flux, pour 156.2 millions de messages par seconde.
\item 449\,500 messages, écoulés en \qty{2.8}{\milli\second}.
\item 64~: la milliseconde la plus chargée contient 609\,500 messages, et 64 cœurs en traitent 640\,000 par milliseconde.
\item Ses 100 millisecondes les plus chargées, 13.56 millions de messages contre les 13.575 millions de l'avis.
\end{enumerate}
\textbf{Partie III.}
\begin{enumerate}[resume]
\item 2.4 millions de cotations, 48\,000 messages de cotation de masse et \qty{0.31}{\giga\bit\per\second}.
\item Un pour le débit (un port transporte 50\,000 messages par seconde)~; douze, un par moteur, selon la structure de la plateforme, et la firme
  utilise les deux ports alloués par moteur, soit 24.
\item \qty{10.0}{\milli\second}, \qty{10.2}{\milli\second} et \qty{15.0}{\milli\second}.
\item Pendant ce délai, certaines de ses cotations sont périmées, et quiconque est plus rapide peut les négocier~; la firme a besoin de purges et
  d'un ordre de rafraîchissement qui retire en premier les cotations les plus dangereuses.
\end{enumerate}
\textbf{Partie IV.}
\begin{enumerate}[resume]
\item \emph{Résultat nommé}~: coter 1\,200 classes de 250 séries à quatre mises à jour par seconde et par côté nécessite 2.4 millions de cotations et 48\,000
  messages de cotation de masse par seconde, ce qu'un seul port pourrait porter, mais que douze moteurs transforment en 24 ports de cotation (deux par moteur), rafraîchissant chaque cotation en 10 à \qty{15}{\milli\second} selon la couche de risque~;
  prendre le flux consolidé dans son intégralité, les deux flux, nécessite 44
  liens de \qty{10}{\giga\bit\per\second} au 99.99\up{e}~centile des millisecondes (12 au 99\up{e}) et 16 cœurs de handler par flux.
\item Publié~: la projection de capacité, la mesure de la milliseconde la plus chargée, le nombre de liens, les règles et tarifs des ports, les
  latences des passerelles. Supposé~: la part de chaque ligne, la forme du modèle de rafales, les octets comme charge utile, la vitesse des cœurs, l'aller-retour, la couche logicielle, la charge de la firme et les
  douze moteurs de la plateforme.
\item La milliseconde la plus chargée demande 32 liens au lieu de 178, et le 99.99\up{e}~centile 16 au lieu de 44~; moins de liens, mais plus rapides, avec les
  mêmes ports de commutateur.
\item Le temps de rafraîchissement complet, donc la durée pendant laquelle ses cotations sont périmées après un mouvement~; la latence d'un message compte moins.
\item Les plus dangereuses d'abord~: les côtés les plus susceptibles d'être touchés (proches de la monnaie, à courte échéance, ceux que le mouvement a rendus bon marché), en purgeant ce qui
  ne peut pas être mis à jour à temps.
\item Le plan de flux (liens, cœurs, file) issu de la nouvelle projection, l'adéquation aux dernières mesures et les tarifs des ports.
\item Dans les lignes des actions et des options~: liens et cœurs pour le flux, ports par plateforme, couche de risque, et leurs coûts mensuels.
\item Parce que les contraintes déterminantes sont les rafales du flux et les allers-retours des ports, qui fixent le nombre de liens, de cœurs et de ports que la
  firme doit acheter avant de gagner la moindre microseconde.
\end{enumerate}
\end{solution}

\begin{solution}{iq:nw:equities-and-options-connectivity:1}
Parce que les récepteurs sont dimensionnés pour les rafales~: une moyenne sur une seconde masque des pics plusieurs fois supérieurs, et un port de commutateur ou un handler qui
 ne suit pas le rythme pendant quelques millisecondes perd ou retarde des données. Même 10 millisecondes masquent les pics à la milliseconde.
\iqlookfor{Moyennes face aux pics~; tampons~; fenêtre de planification.}
\end{solution}

\begin{solution}{iq:nw:equities-and-options-connectivity:2}
Prendre l'avis de capacité et les mesures de l'opérateur, convertir en bits sur le support, doubler pour les deux flux, ajouter la marge de retransmission~;
modéliser les rafales à la milliseconde par \omterm{def:nw:equities-and-options-connectivity:partition}{partition du flux}, choisir un centile de planification, regrouper les partitions sur les liens~; dimensionner les cœurs de handler sur le pic
et vérifier les files~; recalculer à chaque avis.
\iqlookfor{Capacité publiée~; encapsulation~; centiles~; partitions~; handlers.}
\end{solution}

\begin{solution}{iq:nw:equities-and-options-connectivity:3}
Le nombre de cotations à modifier, le nombre de cotations par message, les ports par moteur et l'aller-retour par message (couche de risque comprise),
et les limiteurs de messages de la plateforme. Mesurer en horodatant chaque message de cotation de masse et son accusé de réception par port, ainsi que le délai entre un événement de marché
jusqu'à la mise à jour de la dernière cotation.
\iqlookfor{Messages et allers-retours~; ports~; couche de risque~; mesure de l'événement à la dernière cotation.}
\end{solution}

\begin{solution}{iq:nw:equities-and-options-connectivity:4}
La latence ajoutée dans chaque sens aux centiles utiles, avec et sans charge~; les contrôles qu'il effectue et la manière dont ils sont configurés~;
le comportement en cas de panne~; la capacité par port~; et qui l'exploite. Comparer selon le temps de rafraîchissement complet et l'aller-retour que chacune implique pour la
cotation de la firme, et non seulement par la latence d'un seul message.
\iqlookfor{Distribution de la latence sous charge~; contrôles~; modes de panne~; effet sur le débit.}
\end{solution}

\begin{solution}{iq:nw:equities-and-options-connectivity:5}
Les cotations les plus susceptibles d'être négociées contre des prix périmés~: proches de la monnaie, à courte échéance, de grandes tailles, et les côtés que le mouvement a rendus trop
 généreux~; purger ce qui ne peut pas être mis à jour à temps plutôt que de le laisser.
\iqlookfor{Sélection adverse~; ordre par risque~; purges.}
\end{solution}

\begin{solution}{iq:nw:equities-and-options-connectivity:6}
Les deux flux sur des liens dimensionnés à un centile de millisecondes choisi, avec des tampons de commutateur et une surveillance des pertes~; cœurs de handler dimensionnés avec un
 budget de file~; ports de cotation sur chaque moteur, ports de purge, un drop copy~; une couche de risque mesurée pour son effet sur le temps de rafraîchissement~; et
 surveillance des trous, des délais de file, des temps d'accusé de réception et des limites de ports, replanifiée à chaque avis de capacité.
\iqlookfor{Les deux budgets~; centiles et files~; ports par moteur~; couche de risque~; surveillance.}
\end{solution}
